\documentclass[12pt,a4paper]{report}

%% ─── FONTS & ENCODING ───────────────────────────────────────────────────────
\usepackage[T1]{fontenc}
\usepackage[utf8]{inputenc}
\usepackage{lmodern}
\usepackage{microtype}

%% ─── PAGE GEOMETRY ──────────────────────────────────────────────────────────
\usepackage[a4paper,
            top=2.5cm, bottom=2.5cm,
            left=3cm,  right=2.5cm,
            headheight=15pt]{geometry}

%% ─── COLORS ─────────────────────────────────────────────────────────────────
\usepackage[dvipsnames,svgnames,table]{xcolor}
\definecolor{primary}{RGB}{31,73,125}
\definecolor{accent}{RGB}{0,112,192}
\definecolor{lightgray}{RGB}{245,245,245}
\definecolor{midgray}{RGB}{220,220,220}
\definecolor{darkgray}{RGB}{80,80,80}
\definecolor{codered}{RGB}{192,0,0}
\definecolor{codebg}{RGB}{248,248,248}
\definecolor{screenshotbg}{RGB}{230,238,250}
\definecolor{screenshotborder}{RGB}{31,73,125}
\definecolor{diagrambg}{RGB}{235,245,235}
\definecolor{diagramborder}{RGB}{34,100,34}

%% ─── SPACING & LAYOUT ───────────────────────────────────────────────────────
\usepackage{setspace}
\onehalfspacing
\setlength{\parindent}{1.5em}
\setlength{\parskip}{0.55em}

%% ─── MATH ────────────────────────────────────────────────────────────────────
\usepackage{amsmath,amssymb}

%% ─── GRAPHICS ───────────────────────────────────────────────────────────────
\usepackage{graphicx}
\usepackage{grffile}
\usepackage{float}
\usepackage{caption}
\usepackage{subcaption}
\captionsetup{
  font={small,it},
  labelfont={bf,color=primary},
  labelsep=period,
  skip=8pt
}

%% ─── TABLES ─────────────────────────────────────────────────────────────────
\usepackage{booktabs}
\usepackage{longtable}
\usepackage{array}
\usepackage{tabularx}
\usepackage{multirow}
\renewcommand{\arraystretch}{1.5}

%% ─── LISTINGS / CODE ─────────────────────────────────────────────────────────
\usepackage{listings}
\lstset{
  basicstyle=\ttfamily\small,
  backgroundcolor=\color{codebg},
  frame=single,
  framesep=4pt,
  rulecolor=\color{midgray},
  keywordstyle=\color{accent}\bfseries,
  stringstyle=\color{codered},
  commentstyle=\color{darkgray}\itshape,
  breaklines=true,
  breakatwhitespace=true,
  numbers=left,
  numberstyle=\tiny\color{darkgray},
  xleftmargin=1.5em,
}

%% ─── HEADERS & FOOTERS ───────────────────────────────────────────────────────
\usepackage{fancyhdr}
\pagestyle{fancy}
\fancyhf{}
\fancyhead[L]{\small\color{darkgray}\nouppercase{\leftmark}}
\fancyhead[R]{\small\color{darkgray}MES-X Dependency Management}
\fancyfoot[C]{\small\color{darkgray}\thepage}
\fancyfoot[R]{\small\color{darkgray}ESPRIT -- 2024/2025}
\renewcommand{\headrulewidth}{0.5pt}
\renewcommand{\footrulewidth}{0.5pt}
\renewcommand{\headrule}{\hbox to\headwidth{\color{primary}\leaders\hrule height \headrulewidth\hfill}}
\renewcommand{\footrule}{\hbox to\headwidth{\color{primary}\leaders\hrule height \footrulewidth\hfill}}

%% ─── TOC / MINITOC ───────────────────────────────────────────────────────────
\usepackage{tocloft}
\usepackage{minitoc}
\usepackage{titletoc}
\setcounter{tocdepth}{3}
\setcounter{secnumdepth}{3}
\renewcommand{\cfttoctitlefont}{\Large\bfseries\color{primary}}
\renewcommand{\cftloftitlefont}{\Large\bfseries\color{primary}}
\renewcommand{\cftlottitlefont}{\Large\bfseries\color{primary}}
\renewcommand{\cftchapfont}{\bfseries\color{primary}}
\renewcommand{\cftchappagefont}{\bfseries\color{primary}}
\renewcommand{\contentsname}{Table of Contents}
\renewcommand{\listtablename}{List of Tables}
\renewcommand{\listfigurename}{List of Figures}
\renewcommand{\cftchappresnum}{Chapter\ }
\renewcommand{\cftchapaftersnum}{.\ }
\setlength{\cftchapnumwidth}{7em}
\renewcommand{\mtctitle}{Contents}
\dominitoc

%% ─── CHAPTER / SECTION TITLES ────────────────────────────────────────────────
\usepackage{titlesec}

\titleformat{\chapter}[display]
  {\normalfont\bfseries\color{primary}}
  {\filleft\Large\color{accent}\MakeUppercase{\chaptertitlename}\ 
   \fontsize{60}{60}\selectfont\thechapter}
  {-10pt}
  {\titlerule[1.5pt]\vspace{6pt}\filright\fontsize{18}{22}\selectfont}
  [\vspace{4pt}\titlerule]

\titlespacing*{\chapter}{0pt}{10pt}{30pt}

\titleformat{\section}
  {\normalfont\fontsize{14}{17}\selectfont\bfseries\color{primary}}
  {\thesection}{1em}{}
  [\color{accent}\titlerule[0.4pt]]

\titleformat{\subsection}
  {\normalfont\fontsize{12}{15}\selectfont\bfseries\color{accent}}
  {\thesubsection}{1em}{}

\titleformat{\subsubsection}
  {\normalfont\fontsize{11}{14}\selectfont\bfseries\color{darkgray}}
  {\thesubsubsection}{1em}{}

%% ─── HYPERLINKS ──────────────────────────────────────────────────────────────
\usepackage[hyphens]{url}
\usepackage[
  colorlinks=true,
  linkcolor=primary,
  urlcolor=accent,
  citecolor=accent,
  bookmarks=true,
  bookmarksopen=true,
  pdftitle={MES-X Dependency Management -- Internship Report},
  pdfauthor={Alaa eddine Mdalla},
  pdfsubject={First Internship Report},
  pdfkeywords={NestJS, Azure DevOps, Vue 3, dependency management, ESPRIT}
]{hyperref}

%% ─── MISC PACKAGES ───────────────────────────────────────────────────────────
\usepackage{pdfpages}
\usepackage{ragged2e}
\usepackage{indentfirst}
\usepackage{enumitem}
\setlist[itemize]{leftmargin=2em, itemsep=3pt, topsep=5pt}
\setlist[enumerate]{leftmargin=2em, itemsep=3pt, topsep=5pt}

%% ─── COLORED BOXES ───────────────────────────────────────────────────────────
\usepackage{tcolorbox}
\tcbuselibrary{skins,breakable}

%% Info box
\newtcolorbox{infobox}[1][]{
  enhanced, breakable,
  colback=screenshotbg, colframe=primary,
  boxrule=0.8pt, arc=3pt,
  left=10pt, right=10pt, top=8pt, bottom=8pt,
  fonttitle=\bfseries\small\color{white},
  colbacktitle=primary,
  title=#1
}

%% Diagram placeholder box (use-case or class diagrams)
\newcommand{\umlplaceholder}[2][UML Diagram Placeholder]{%
  \begin{figure}[H]
    \centering
    \begin{tcolorbox}[
      enhanced,
      colback=diagrambg,
      colframe=diagramborder,
      boxrule=1.2pt,
      arc=5pt,
      left=6pt, right=6pt, top=6pt, bottom=6pt,
      width=\textwidth,
      height=8cm,
      valign=center,
    ]
      \centering
      \color{diagramborder}
      \large\bfseries [\,#1\,]\\[8pt]
      \normalsize\normalfont\color{darkgray}
      Insert your diagram image here:\\[4pt]
      \ttfamily\small\color{darkgray}\textbackslash includegraphics[width=\textbackslash textwidth]\{img/#2\}
    \end{tcolorbox}
    \caption{#1}
    \label{fig:#2}
  \end{figure}%
}

%% Screenshot placeholder
\newcommand{\screenshot}[2][Screenshot placeholder]{%
  \begin{figure}[H]
    \centering
    \begin{tcolorbox}[
      enhanced,
      colback=screenshotbg,
      colframe=screenshotborder,
      boxrule=1pt,
      arc=4pt,
      left=4pt, right=4pt, top=4pt, bottom=4pt,
      width=\textwidth,
      height=7cm,
      valign=center,
    ]
      \centering
      \color{primary}
      \large\bfseries [\,SCREENSHOT\,]\\[6pt]
      \normalsize\normalfont\color{darkgray}#1
    \end{tcolorbox}
    \caption{#2}
  \end{figure}%
}

%% Wide screenshot placeholder
\newcommand{\screenshotwide}[2][Screenshot placeholder]{%
  \begin{figure}[H]
    \centering
    \begin{tcolorbox}[
      enhanced,
      colback=screenshotbg,
      colframe=screenshotborder,
      boxrule=1pt,
      arc=4pt,
      left=4pt, right=4pt, top=4pt, bottom=4pt,
      width=\textwidth,
      height=9cm,
      valign=center,
    ]
      \centering
      \color{primary}
      \large\bfseries [\,SCREENSHOT\,]\\[6pt]
      \normalsize\normalfont\color{darkgray}#1
    \end{tcolorbox}
    \caption{#2}
  \end{figure}%
}

%% Inline code
\newcommand{\code}[1]{\texttt{\small\color{codered}#1}}

%% ═══════════════════════════════════════════════════════════════════════════
%%  DOCUMENT
%% ═══════════════════════════════════════════════════════════════════════════
\begin{document}

\begin{titlepage}
\thispagestyle{empty}

% Remove margins ONLY for this page
\newgeometry{margin=0pt}

\noindent

\includegraphics[width=\paperwidth,height=\paperheight]{img/cover/cover.png}

% Restore normal margins
\restoregeometry

\end{titlepage}

%% ─── FRONT MATTER ────────────────────────────────────────────────────────────
\pagenumbering{roman}
\pagestyle{fancy}

%% ── DEDICATION ──────────────────────────────────────────────────────────────
\newpage
\thispagestyle{empty}
\vspace*{3cm}
\begin{center}
  {\Large\bfseries\color{primary} Dedication}\\[0.5cm]
  \textcolor{accent}{\rule{0.3\textwidth}{1pt}}
\end{center}
\vspace{2cm}

\begin{center}
\begin{minipage}{0.82\textwidth}
\centering\itshape\large\setstretch{1.8}

I dedicate this work to my parents, who have always supported me throughout my
studies and given me the confidence to keep moving forward.

\medskip
To my father - for his strength, discipline, and constant encouragement.

To my mother - for her patience, love, and unwavering prayers.

\medskip
To my sister Wafa and my brother Dhia, for their kindness and steady support.

\medskip
To everyone who believed in me and helped me stay focused until the end of this journey.

\bigskip
\normalfont\normalsize\color{darkgray}
And finally, to myself --- because I reached this stage through persistence,
discipline, and hard work.

\end{minipage}
\end{center}

%% ── ACKNOWLEDGEMENTS ────────────────────────────────────────────────────────
\newpage
\thispagestyle{empty}
\vspace*{3cm}
\begin{center}
  {\Large\bfseries\color{primary} Acknowledgements}\\[0.5cm]
  \textcolor{accent}{\rule{0.3\textwidth}{1pt}}
\end{center}
\vspace{1.8cm}

\begin{minipage}{0.88\textwidth}
\setstretch{1.7}

I would like to express my sincere gratitude to everyone who supported and guided
me during this internship.

\medskip
First, I thank my academic supervisor \textbf{Wafa Boumaiza} for her thoughtful
feedback, consistent availability, and the guidance that helped me continually
improve the quality of this work.

\medskip
I also extend my thanks to \textbf{ESPRIT} for the solid academic foundation that
prepared me to take on a real-world engineering challenge.

\medskip
I am grateful to the members of the project and development teams for their
technical input, constructive discussions, and the collaborative environment they
created throughout the internship period.

\medskip
Most importantly, I thank the people closest to me for their encouragement,
patience, and support during this demanding but rewarding experience.

\end{minipage}

\vspace{1.5cm}
\begin{flushright}
\itshape Alaa eddine Mdalla
\end{flushright}

%% ── TABLE OF CONTENTS ───────────────────────────────────────────────────────
\newpage
\tableofcontents

\newpage
\addcontentsline{toc}{chapter}{List of Figures}
\listoffigures

\newpage
\addcontentsline{toc}{chapter}{List of Tables}
\listoftables

\newpage

%% ── LIST OF ABBREVIATIONS ───────────────────────────────────────────────────
\chapter*{List of Abbreviations}
\addcontentsline{toc}{chapter}{List of Abbreviations}
\markboth{List of Abbreviations}{List of Abbreviations}

\vspace{0.5cm}
\begin{center}
\begin{tabular}{@{}>{\bfseries\color{primary}}p{2.8cm} p{11.5cm}@{}}
\toprule
\textbf{\color{primary}Abbreviation} & \textbf{\color{primary}Definition} \\
\midrule
API         & Application Programming Interface \\
Azure DevOps & Microsoft platform for source control, work tracking, and CI/CD \\
CI/CD       & Continuous Integration / Continuous Delivery \\
CORS        & Cross-Origin Resource Sharing \\
Dapr        & Distributed Application Runtime \\
DTO         & Data Transfer Object \\
ESPRIT      & Higher School of Engineering and Technologies \\
GET         & HTTP method used to retrieve resources \\
Jest        & JavaScript testing framework used throughout the project \\
MES-X       & The platform context of this dependency-management service \\
NestJS      & Node.js framework used to build the backend \\
OIDC        & OpenID Connect \\
PAT         & Personal Access Token \\
PR          & Pull Request \\
SSE         & Server-Sent Events \\
Swagger     & OpenAPI documentation interface generated by the backend \\
TypeScript  & Statically typed superset of JavaScript used for the implementation \\
UI          & User Interface \\
UML         & Unified Modeling Language \\
WIQL        & Work Item Query Language \\
\bottomrule
\end{tabular}
\end{center}

%% ═══════════════════════════════════════════════════════════════════════════
%%  MAIN MATTER
%% ═══════════════════════════════════════════════════════════════════════════
\newpage
\pagenumbering{arabic}

%% ════════════════════════════════════════════════════════════════════════════
\chapter*{General Introduction}
\addcontentsline{toc}{chapter}{General Introduction}
\markboth{General Introduction}{General Introduction}

\vspace{0.4cm}

Modern software teams rely heavily on structured traceability to understand how
work items, pull requests, commits, branches, and tags relate to one another.
Without a dedicated service, this investigation is manual, repetitive, and
error-prone. The \textbf{MES-X Dependency Management} project was born out of
exactly this need.

\medskip
The solution spans two repositories within the same workspace: the NestJS backend
in \code{mesx-proto-dependency-mgt-be} and the Vue\,3 trace UI in
\code{mesx-proto-dependency-mgt-ui}. The backend serves as the authoritative
integration layer with Azure DevOps, while the UI exposes the traceability
features through dedicated screens for single work items, batch task loading, and
batch related-item analysis.

\medskip
This was not a classic CRUD project. The objective was to build a read-only
dependency-management service capable of answering practical questions: which
commits are linked to a work item, which work items are linked to a pull request,
which branches and tags belong to a repository, how a work-item graph can be
traversed recursively, and how results can be rendered progressively in the
frontend.

\bigskip
\noindent This report is organized into ten chapters:

\begin{itemize}
  \item \textbf{Chapter 1 - General Presentation:} project context, objectives, repositories, and delivered scope.
  \item \textbf{Chapter 2 - Requirements and Architecture:} functional and non-functional constraints, module structure, and cross-repository flow.
  \item \textbf{Chapter 3 - Work Item Module:} endpoint design, recursive traversal, SSE streaming, and frontend usage.
  \item \textbf{Chapter 4 - Pull Request Module:} PR lookup, linked work-item resolution, and repository-summary enrichment.
  \item \textbf{Chapter 5 - Commit Module:} commit aggregation, branch and tag discovery, and batch tag resolution.
  \item \textbf{Chapter 6 - Decision Module:} aggregation input contract, classification rules, and repository-level decision output.
  \item \textbf{Chapter 7 - Azure DevOps Integration:} secret retrieval, configuration modules, Swagger, validation, and CORS.
  \item \textbf{Chapter 8 - Performance, Caching, and Reliability:} batching, concurrency limits, deduplication, and error handling.
  \item \textbf{Chapter 9 - Testing, Tooling, and Delivery:} tests, scripts, Docker, pipelines, and frontend integration.
  \item \textbf{Chapter 10 - Final Result and Conclusion:} outcome, lessons learned, limitations, and future improvements.
\end{itemize}

\section*{General Framework}

This project was completed as part of engineering studies at \textbf{ESPRIT}. It
offered the opportunity to work on a real dependency-management problem, deepen
practical knowledge of Azure DevOps APIs, structure a maintainable NestJS backend,
and align that backend with a Vue-based trace UI.

%% ════════════════════════════════════════════════════════════════════════════
\chapter{General Presentation of the Project}
\startcontents
\printcontents{}{}{}
\newpage

\section{Introduction}

This chapter introduces the MES-X Dependency Management project, the repositories
involved, the business need it addresses, and the concrete scope of the delivered
solution.

%% ── CONTEXT ────────────────────────────────────────────────────────────────
\section{Project Context}

Software teams that use Azure DevOps need to understand relationships between work
items, pull requests, commits, branches, and tags. Without a dedicated integration
layer, tracing these relationships requires navigating multiple Azure DevOps screens
manually - a time-consuming and error-prone process.

\medskip
The MES-X Dependency Management solution centralizes these relationships through a
consistent backend API and surfaces them through an operational trace UI with three
main screens:

\begin{itemize}
  \item a single work-item trace screen;
  \item a batch task-loading screen driven by SSE streaming;
  \item a batch related-results screen that aggregates graphs and enriches commits
        with tag context.
\end{itemize}

\begin{figure}[H]
  \centering
  \includegraphics[width=\textwidth]{img/MES.png}
  \caption{General platform overview - MES-X Dependency Management service in context}
\end{figure}

%% ── REPOSITORIES ────────────────────────────────────────────────────────────
\section{Repositories and Delivered Scope}

The delivered solution is distributed across two repositories, as summarized in
Table~\ref{tab:workspace-repositories}.

\begin{table}[H]
\centering
\caption{Workspace repositories and their roles}
\label{tab:workspace-repositories}
\rowcolors{2}{lightgray}{white}
\begin{tabularx}{\textwidth}{|>{\ttfamily\small}p{5.2cm}|p{2.8cm}|X|}
\hline
\rowcolor{primary}
  \textcolor{white}{\textbf{Repository}} &
  \textcolor{white}{\textbf{Stack}} &
  \textcolor{white}{\textbf{Role}} \\
\hline
mesx-proto-dependency-mgt-be &
  NestJS 10, TypeScript, Swagger, Dapr &
  Core backend service exposing traceability endpoints and Azure DevOps integration \\
mesx-proto-dependency-mgt-ui &
  Vue\,3, Quasar\,2, Pinia, Vite, Axios &
  Trace UI exposing single-item, batch-stream, and batch-related analysis views \\
\hline
\end{tabularx}
\end{table}

%% ── USE-CASE DIAGRAM ────────────────────────────────────────────────────────
\section{Use-Case Diagram}

The following diagram illustrates the main actors and use cases of the system.
The two primary actors are the \textbf{End User} (a developer or team member using
the trace UI) and the \textbf{Azure DevOps Platform} (the external data source).

\begin{figure}[p]
    \centering
    \vspace*{\fill}
    \includegraphics[width=0.95\textwidth, height=0.92\textheight, keepaspectratio]{img/archi/use_case.png}
    \caption{Use Case Diagram - MES-X Dependency Management}
    \label{fig:usecase_diagram}
    \vspace*{\fill}
\end{figure}
\noindent The main use cases include:
\begin{itemize}
  \item \textit{Trace a single work item} - visualize its related graph and linked commits;
  \item \textit{Load task items via batch stream} - progressively receive work items through SSE;
  \item \textit{Analyze batch related results} - select multiple roots and explore aggregated dependency data;
  \item \textit{Resolve commit tags} - find the exact or surrounding tags for one or a batch of commits.
\end{itemize}

%% ── MY ROLE ────────────────────────────────────────────────────────────────
\section{My Role in the Project}

My contribution covered both the backend and the frontend alignment. On the backend
side, I structured the domain modules, implemented controllers and services,
normalized Azure DevOps payloads into stable DTOs, added traversal and batch
algorithms, and set up the shared configuration and logging layers. On the frontend
side, I aligned the three trace screens with the backend contracts and documented
the real integration behavior between the two repositories.

%% ── EXISTING NEED ───────────────────────────────────────────────────────────
\section{Study of the Existing Need}

Before building the solution, I identified the following concrete needs:

\begin{itemize}
  \item retrieve work items by ID with normalized relation data;
  \item navigate recursively from a root work item to its children and PR-linked items;
  \item stream large task result sets progressively rather than blocking the UI;
  \item resolve work items linked to a given pull request;
  \item retrieve commits linked to work items or pull requests;
  \item list repository branches and tags;
  \item resolve exact or surrounding tags for one or more commits;
  \item expose all of the above through a documented API and a usable trace UI.
\end{itemize}

%% ── PROBLEM STATEMENT ───────────────────────────────────────────────────────
\section{Problem Statement}

The core challenge is that relationship data in Azure DevOps is scattered across
several entity types and APIs. A work item may link directly to commits, indirectly
to pull requests, and recursively to child work items. The UI needs deterministic
and structured responses, not raw Azure DevOps payloads.

\medskip
The solution therefore had to:
\begin{itemize}
  \item connect to Azure DevOps securely and reliably;
  \item normalize heterogeneous data into stable DTO contracts;
  \item avoid duplicating work items, commits, and pull-request traversals;
  \item remain maintainable and testable as features were added;
  \item support both single-item and batch-oriented user flows.
\end{itemize}

%% ── PROPOSED SOLUTION ───────────────────────────────────────────────────────
\section{Proposed Solution}

To address these needs, I implemented a modular backend with four domain modules
and an integrated trace UI:

\begin{itemize}
  \item \textbf{Work Item Module} - work-item lookup, recursive related-item traversal, and SSE task streaming;
  \item \textbf{Pull Request Module} - PR-to-work-item resolution and repository-summary enrichment;
  \item \textbf{Commit Module} - commit aggregation, branch and tag discovery, single and batch tag resolution;
  \item \textbf{Decision Module} - post-processing aggregation that classifies work items and repository commits for decision support;
  \item \textbf{Trace UI} - three dedicated routes for single, batch-stream, and batch-related analysis.
\end{itemize}

%% ── DEVELOPMENT APPROACH ────────────────────────────────────────────────────
\section{Development Approach}

I followed an incremental approach: first securing the Azure DevOps connection
and configuration layer, then building the work-item flows, then pull-request
resolution, then commit and tag features, then implementing the decision
aggregation layer, and finally aligning the frontend trace flow with the actual
backend behavior.

\subsection{Methodological Choice}

An iterative workflow suited this project better than a one-shot implementation
because the modules are interdependent. Commit aggregation reuses work-item and
PR logic; the work-item traversal depends on relation parsing, PR expansion, cache
behavior, and commit aggregation. The Decision module then depends on structured
inputs produced by prior modules (work-item state sets and repository tag
position data), which reinforced the need for progressive validation. Developing
feature by feature allowed issues to be caught and corrected early.

\begin{table}[H]
\centering
\caption{Development approach comparison}
\label{tab:dev-approach}
\rowcolors{2}{lightgray}{white}
\begin{tabular}{|p{3.2cm}|p{5.6cm}|p{5.6cm}|}
\hline
\rowcolor{primary}
  \textcolor{white}{\textbf{Criteria}} &
  \textcolor{white}{\textbf{Iterative approach}} &
  \textcolor{white}{\textbf{One-shot approach}} \\
\hline
Change handling    & Easy to adapt while building and validating         & Hard to correct after the fact \\
Validation         & Tested feature by feature                           & Validation happens late \\
Risk               & Lower - issues are isolated earlier                 & Higher - problems accumulate \\
Cross-module align & Backend and UI synchronized incrementally           & More likely to produce contract drift \\
Maintainability    & Encourages clean abstractions                       & Encourages technical debt \\
\hline
\end{tabular}
\end{table}

\section{Conclusion}

This chapter presented the context, repositories, scope, and the problem addressed
by the project. The next chapter translates this need into concrete technical
requirements and explains the implemented architecture.

%% ════════════════════════════════════════════════════════════════════════════
\chapter{Requirements Analysis and Architecture}

\section{Introduction}

This chapter translates the identified needs into concrete requirements and
describes the architecture implemented across the backend and the trace UI.

%% ── FUNCTIONAL REQUIREMENTS ─────────────────────────────────────────────────
\section{Functional Requirements}

The solution had to provide the following read-only capabilities:

\begin{itemize}
  \item get a work item by ID;
  \item traverse related work items recursively for one root;
  \item traverse multiple roots and return one aggregated response;
  \item stream task work items progressively through SSE;
  \item get work items linked to a pull request;
  \item get commits linked to a work item or pull request;
  \item list repository branches and tags;
  \item resolve exact or surrounding tags for a single commit;
  \item resolve exact or surrounding tags for a batch of up to 100 commits;
  \item aggregate decision input into done/other work-item sets and repository commit groups;
  \item expose all of the above through the trace UI.
\end{itemize}

%% ── NON-FUNCTIONAL REQUIREMENTS ─────────────────────────────────────────────
\section{Non-Functional Requirements}

Beyond feature delivery, the implementation had to respect the following
technical constraints:

\begin{itemize}
  \item clean separation between controllers, services, DTOs, and helpers;
  \item safe secret retrieval and external configuration loading via Dapr;
  \item bounded concurrency when expanding pull requests or resolving commits;
  \item deterministic, deduplicated result sets across all endpoints;
  \item explicit input validation and structured error responses;
  \item observability through a centralized logging layer;
  \item testability, documented API contracts, and reproducible delivery scripts.
\end{itemize}

%% ── ARCHITECTURE OVERVIEW ───────────────────────────────────────────────────
\section{Architecture Overview}

The backend is organized around \code{AppModule}, which imports the following
major sub-modules:

\begin{center}
\begin{tcolorbox}[
  enhanced, colback=screenshotbg, colframe=primary, boxrule=0.8pt, arc=4pt,
  width=0.80\textwidth, left=12pt, right=12pt, top=10pt, bottom=10pt,
  title={\textbf{AppModule - imported modules}},
  fonttitle=\bfseries\small\color{white},
  colbacktitle=primary
]
\begin{itemize}[itemsep=2pt]
  \item \code{WorkItemModule}
  \item \code{PullRequestModule}
  \item \code{CommitModule}
  \item \code{DecisionModule}
  \item \code{AppLoggerModule}
  \item \code{AppConfigModule}
  \item \code{AzureConfigModule}
  \item \code{UrlApiConfigModule}
\end{itemize}
\end{tcolorbox}
\end{center}

\medskip
The frontend uses Vue Router, Pinia, a shared Axios-based API layer, and
dedicated composables for the batch trace flow.

\begin{figure}[H]
  \centering
  \includegraphics[width=\textwidth]{img/archi/diagramm.png}
  \caption{High-level architecture: NestJS backend and Vue 3 trace UI}
    \label{fig::High-level architecture}
\end{figure}

%% ── CLASS DIAGRAM ───────────────────────────────────────────────────────────
\section{Class Diagram}

The class diagram below shows the main domain classes and services, their
attributes, and the relationships between them. The four domain modules
(\textit{WorkItem}, \textit{PullRequest}, \textit{Commit}, \textit{Decision}) each follow the same
layered pattern: \textbf{Controller} → \textbf{Service} → \textbf{DTO}.

\begin{figure}[H]
  \centering
  \includegraphics[width=\textwidth]{img/archi/classdiag.png}
  \caption{Class Diagram}
\end{figure}

\noindent Key relationships visible in the diagram include:
\begin{itemize}
  \item \code{WorkItemService} depends on \code{PullRequestService} for PR expansion during recursive traversal;
  \item \code{CommitWorkItemService} calls \code{CommitPRService} to aggregate commits from PR links;
  \item \code{AzureDevopsService} is injected into all three domain services as the single Azure DevOps connection point;
  \item \code{CacheManagerService} is shared between \code{WorkItemService} and the traversal helpers.
\end{itemize}

%% ── ENDPOINT INVENTORY ──────────────────────────────────────────────────────
\section{Implemented Endpoint Inventory}

Table~\ref{tab:endpoint-inventory} summarizes the full public endpoint surface
of the backend.

\begin{table}[H]
\centering
\caption{Implemented backend endpoint inventory}
\label{tab:endpoint-inventory}
\rowcolors{2}{lightgray}{white}
\begin{tabularx}{\textwidth}{|p{3.2cm}|p{5.5cm}|X|}
\hline
\rowcolor{primary}
  \textcolor{white}{\textbf{Module}} &
  \textcolor{white}{\textbf{Endpoint}} &
  \textcolor{white}{\textbf{Purpose}} \\
\hline
Work Item    & \code{GET /workitem/:id}                     & Get one work item with normalized relations \\
Work Item    & \code{GET /workitem/:id/related}             & Recursive related graph for one root \\
Work Item    & \code{GET /workitem/related}                 & Aggregated graph for multiple roots \\
Work Item    & \code{GET /workitem/tasks/stream}            & Stream task batches progressively over SSE \\
Pull Request & \code{GET /pullrequests/:prId/workitems}     & Resolve work items linked to a PR \\
Commit       & \code{GET /commits/workitems/:workItemId}    & Commits linked to a work item \\
Commit       & \code{GET /commits/pullrequests/:prId}       & Commits linked to a pull request \\
Commit       & \code{GET /commits/branches}                 & List repository branches and tags \\
Commit       & \code{GET /commits/tags}                     & Exact or surrounding tags for one commit \\
Commit       & \code{GET /commits/tags/batch}               & Batch tag resolution for up to 100 commits \\
Decision     & \code{POST /decision/aggregate}              & Aggregate work-item states and repository tag-position groups \\
\hline
\end{tabularx}
\end{table}

%% ── CROSS-REPO FLOW ─────────────────────────────────────────────────────────
\section{Cross-Repository Request Flow}

A typical end-to-end request follows these steps:

\begin{enumerate}
  \item the user opens one of the three trace screens in the UI;
  \item Vue Router activates the corresponding view component;
  \item the view or composable delegates state management to \code{traceStore};
  \item the store calls \code{trace-api.service.ts}, which uses Axios with \code{withCredentials: true};
  \item the backend controller validates and routes the request to the appropriate service;
  \item the domain service calls Azure DevOps through \code{AzureDevopsService};
  \item the raw response is normalized into DTOs and returned to the UI.
\end{enumerate}

%% ── CURRENT UI FLOW ─────────────────────────────────────────────────────────
\section{Current UI Flow}

The three frontend screens map directly to backend capabilities:

\begin{itemize}
  \item \textbf{RelatedWorkItemTrace} - fetches one related-work-item graph via \code{GET /workitem/:id/related};
  \item \textbf{BatchWorkItemTrace} - opens an SSE stream on \code{GET /workitem/tasks/stream} and lets the user progressively select work items;
  \item \textbf{BatchRelatedWorkItemTrace} - reads selected IDs, calls the single related endpoint once per root, groups commits by repository, and enriches them via \code{GET /commits/tags/batch}.
\end{itemize}

\begin{infobox}[Known Gap]
The backend already exposes \code{GET /workitem/related} for aggregated batch traversal, but the current UI batch results page still fans out per-root calls. This is a documented gap planned for a future alignment pass.
\end{infobox}

\section{Conclusion}

The requirements led to a modular, contract-driven architecture where the backend
is the single source of truth and the trace UI consumes stable DTOs on top of it.
The next chapter details the most graph-oriented component: the Work Item module.

%% ════════════════════════════════════════════════════════════════════════════
\chapter{Work Item Module}

\section{Introduction}

The Work Item module is the core graph-navigation component of the backend. It
handles direct work-item retrieval, recursive related-item traversal, task
discovery, progressive SSE streaming, and the aggregated batch endpoint for
multiple root IDs.

%% ── ENDPOINTS ───────────────────────────────────────────────────────────────
\section{Implemented Endpoints}

\begin{table}[H]
\centering
\caption{Work Item module endpoints}
\label{tab:work-item-endpoints}
\rowcolors{2}{lightgray}{white}
\begin{tabularx}{\textwidth}{|p{5cm}|p{1.8cm}|X|}
\hline
\rowcolor{primary}
  \textcolor{white}{\textbf{Endpoint}} &
  \textcolor{white}{\textbf{Method}} &
  \textcolor{white}{\textbf{Purpose}} \\
\hline
\code{/workitem/tasks/stream} & GET & Stream task-like work items progressively over SSE \\
\code{/workitem/:id}          & GET & Return one normalized work item \\
\code{/workitem/:id/related}  & GET & Traverse related items for one root ID \\
\code{/workitem/related}      & GET & Traverse multiple roots, return one aggregated response \\
\hline
\end{tabularx}
\end{table}

%% ── TASK DISCOVERY ──────────────────────────────────────────────────────────
\section{Task Discovery and Streaming}

The streaming flow starts with \code{WorkItemService.getTaskIds(project)}, which
builds a WIQL query selecting work items of type \code{TASK}, \code{Bug}, or
\code{User Story}, ordered by changed date descending.

\medskip
Once IDs are known, \code{StreamManagerService.handleStreamRequest(...)} does the
following:

\begin{enumerate}
  \item sets SSE headers through \code{StreamService.setSseHeaders(...)};
  \item registers the client response object;
  \item reuses an active stream if one is already running;
  \item fetches task IDs for the fixed project value;
  \item starts batched work-item loading through \code{getWorkItems(ids, 200, onBatch)};
  \item broadcasts each batch as \code{\{\ batch,\ batchNumber,\ totalLoaded\ \}};
  \item broadcasts completion as \code{\{\ complete:\ true,\ totalTasks,\ batchesProcessed\ \}}.
\end{enumerate}

\begin{figure}[H]
  \centering
  \includegraphics[width=\textwidth]{img/Batchtaskstreaming screen.png}
  \caption{Batch task loading screen (\texttt{/cnerr-ui/trace/workitem-batch})}
\end{figure}

\noindent The UI consumes this stream through the browser's native \code{EventSource}
API. The Pinia store annotates each received task with the stream batch number and
reception timestamp before rendering it.

%% ── BATCHED RETRIEVAL ───────────────────────────────────────────────────────
\section{Batched Work-Item Retrieval}

The generic fetcher \code{WorkItemService.getWorkItems(...)} has these main
characteristics:

\begin{itemize}
  \item a default batch size of 200 IDs per request;
  \item concurrent batch execution using \code{Promise.all};
  \item one automatic retry per failing batch before throwing an internal error;
  \item optional batch callback support, used by the SSE stream;
  \item normalization of raw Azure DevOps fields into the local \code{WorkItem} DTO.
\end{itemize}

%% ── SINGLE RETRIEVAL ────────────────────────────────────────────────────────
\section{Single Work-Item Retrieval}

The method \code{getWorkItemById(...)} loads one work item and converts it into
the local DTO structure. It uses an in-memory cache when allowed by the caller.
The result includes: \code{id}, \code{title}, \code{type}, \code{state},
\code{assignedTo}, and normalized \code{links}.

\begin{figure}[H]
  \centering
  \includegraphics[width=\textwidth]{img/ui/Capture d'écran 2026-04-29 122558.png}
  \caption{Single work-item trace screen (\texttt{/cnerr-ui/trace/workitem-related})}
\end{figure}

%% ── RELATION PARSING ────────────────────────────────────────────────────────
\section{Relation Parsing}

The shared helper \code{extractLinks(...)} parses Azure DevOps relations and
transforms them into a stable internal structure. PR links are identified by
checking whether the relation is an \code{ArtifactLink} whose URL contains
\code{PullRequestId}, and the PR ID is extracted from the decoded relation URL.

%% ── RECURSIVE TRAVERSAL ─────────────────────────────────────────────────────
\section{Recursive Related-Item Traversal}

The recursive graph traversal is the most important behavior of the module.
The implemented algorithm:

\begin{enumerate}
  \item normalizes the root work-item IDs and removes duplicates;
  \item loads the root items and fails early if one does not exist;
  \item initializes a discovered-ID set, a fetched-work-item cache, a PR-to-work-item cache, and a visited-PR set;
  \item traverses one level at a time through \code{traverseLevel(...)};
  \item fetches only the IDs not already present in the fetched-work-item map;
  \item extracts child work-item IDs from hierarchy links;
  \item extracts PR IDs, marks them as visited, and expands them through the Pull Request module;
  \item computes the next level of IDs while excluding already-discovered ones;
  \item recurses until no new IDs remain;
  \item aggregates commits after traversal is complete.
\end{enumerate}

%% ── PR EXPANSION ────────────────────────────────────────────────────────────
\section{PR Expansion and Concurrency}

Pull-request expansion is performed by \code{expandPullRequests(...)}. The
implementation uses \code{p-limit(5)} to cap concurrent PR expansions. Each PR
ID is resolved to work-item IDs through
\code{PullRequestService.getWorkItemIdsByPR(...)}. Results are cached in a
PR-to-work-item map and reused on later traversal steps.

%% ── SINGLE VS BATCH ─────────────────────────────────────────────────────────
\section{Single vs.\ Batch Related Endpoints}

Both response shapes share the same core payload: \code{visitedPullRequestIds},
\code{visitedPullRequests}, \code{commits}, and \code{commitCount}.

\begin{itemize}
  \item \code{GET /workitem/:id/related} returns one \code{RelatedWorkItemsResponseDto} with a single \code{workItemId};
  \item \code{GET /workitem/related} returns one \code{RelatedWorkItemsBatchResponseDto} with \code{requestedWorkItemIds} and a deduplicated payload across all roots.
\end{itemize}

\screenshot[Batch related-results screen - work items grouped with commits and tag enrichment per repository]{Batch related results screen (\texttt{/cnerr-ui/trace/workitem-batch-related})}

\section{Conclusion}

The Work Item module is the backbone of the project. It provides the recursive
traversal logic, task streaming behavior, and the graph-oriented contracts that
drive the rest of the system.

%% ════════════════════════════════════════════════════════════════════════════
\chapter{Pull Request Module}

\section{Introduction}

This chapter explains how the Pull Request module resolves work items from
Azure DevOps pull requests and enriches the work-item traversal with repository
context.

\section{Implemented Endpoint}

\begin{center}
\code{GET /pullrequests/:prId/workitems}
\end{center}

\section{Core Responsibilities}

The Pull Request module fulfills two main responsibilities:

\begin{itemize}
  \item resolve the work items linked to a given pull request;
  \item enrich PR IDs with repository names through \code{getPullRequestRepositorySummaries(...)}.
\end{itemize}

\section{PR Lookup Strategy}

The helper \code{fetchPRById(...)} uses a fallback-aware lookup strategy:

\begin{enumerate}
  \item attempt \code{gitApi.getPullRequestById(...)} as the primary path;
  \item extract the repository ID from the result;
  \item if the primary call fails, fall back to \code{gitApi.getPullRequest(repoId, pullRequestId, project)};
  \item log warnings when fallbacks fail rather than crashing the whole flow.
\end{enumerate}

\section{Resolving Work Items Linked to a PR}

The workflow for \code{getWorkItemsByPR(...)} is:

\begin{enumerate}
  \item validate that the PR ID is numeric and positive;
  \item resolve the repository and PR metadata;
  \item retrieve raw work-item references using \code{getPullRequestWorkItemRefs(...)};
  \item remove duplicates from the returned ID list;
  \item fetch each work item with \code{WorkItemExpand.Relations};
  \item map results into \code{WorkItemResultDto} objects containing ID, title, and normalized links.
\end{enumerate}

\screenshot[Swagger UI showing the Pull Request endpoint - \texttt{GET /pullrequests/:prId/workitems} with example response]{Swagger documentation for the Pull Request module}

\section{Response Construction}

The public builder \code{getWorkItemsResponseDataByPR(...)} wraps linked work
items in a stable response that includes: the PR ID, the resolved project name,
the work-item count, and the list of linked work items.

\section{PR Repository Summaries}

The method \code{getPullRequestRepositorySummaries(...)} enriches visited PR IDs
with their repository names after traversal. The behavior is best-effort:

\begin{itemize}
  \item when the PR resolves, the real repository name is returned;
  \item when the PR cannot be resolved, the backend returns \code{unknown-repository};
  \item this allows the UI to display meaningful names when available and fall back to PR IDs when not.
\end{itemize}

\section{Conclusion}

The Pull Request module bridges Git-based pull requests and work-item
relationships, and provides the repository-summary enrichment now included in all
related-item trace responses.

%% ════════════════════════════════════════════════════════════════════════════
\chapter{Commit Module}

\section{Introduction}

The Commit module completes the dependency-management story by linking work
items and pull requests back to concrete code changes and repository references.

\section{Implemented Endpoints}

\begin{table}[H]
\centering
\caption{Commit module endpoints}
\label{tab:commit-endpoints}
\rowcolors{2}{lightgray}{white}
\begin{tabularx}{\textwidth}{|p{5cm}|p{1.8cm}|X|}
\hline
\rowcolor{primary}
  \textcolor{white}{\textbf{Endpoint}} &
  \textcolor{white}{\textbf{Method}} &
  \textcolor{white}{\textbf{Purpose}} \\
\hline
\code{/commits/workitems/:workItemId} & GET & Commits linked to a work item \\
\code{/commits/pullrequests/:prId}    & GET & Commits linked to a pull request \\
\code{/commits/branches}              & GET & Branches and tags for one repository \\
\code{/commits/tags}                  & GET & Exact or surrounding tags for one commit \\
\code{/commits/tags/batch}            & GET & Batch tag resolution for up to 100 commits \\
\hline
\end{tabularx}
\end{table}

\section{Service Decomposition}

The module is split into focused services to keep each responsibility isolated:

\begin{itemize}
  \item \code{CommitService} - orchestration and input validation;
  \item \code{CommitPRService} - PR-to-commit resolution;
  \item \code{CommitWorkItemService} - work-item-to-commit aggregation;
  \item \code{CommitHistoryService} - branch/tag discovery and tag resolution algorithms.
\end{itemize}

\section{Commits Linked to a Work Item}

\code{CommitWorkItemService.getCommitsByWorkitem(...)} works in two phases:

\begin{enumerate}
  \item load the work item and inspect its normalized links;
  \item collect direct commit links and pull-request IDs separately.
\end{enumerate}

For PR links, the service expands each PR through \code{CommitPRService} using
\code{p-limit(5)} to cap concurrency. The final response is deduplicated through
a \code{Map<string, CommitInfo>} keyed by commit SHA.

\screenshot[Commit list view - commits associated with a work item, showing SHA, message, and linked PR]{Commits linked to a work item}

\section{Commits Linked to a Pull Request}

The endpoint \code{GET /commits/pullrequests/:prId} delegates to the PR commit
service and returns a stable \code{CommitsByPRResponseDto}. It is used both
directly by the UI for PR-centered traces and indirectly inside work-item commit
aggregation.

\section{Repository Branch and Tag Listing}

The endpoint \code{GET /commits/branches} resolves the repository ID, fetches
all references, separates branches from tags, strips the \code{refs/heads/} and
\code{refs/tags/} prefixes, and returns a \code{RepoReferencesResponseDto}.

\section{Single Commit Tag Resolution}

The algorithm inside \code{CommitHistoryService.getBranchesAndTagsByCommit(...)}:

\begin{enumerate}
  \item resolve the repository ID and fetch all refs;
  \item split refs into branches and tags;
  \item verify the target commit exists in the repository;
  \item resolve lightweight and annotated tag references to actual commit IDs;
  \item build a commit-to-tag map;
  \item scan the \code{main} branch first;
  \item if the commit is on \code{main}, return exact or surrounding tags immediately;
  \item otherwise, scan the remaining containing branches;
  \item if no tags are found anywhere, return the containing branch names as a fallback.
\end{enumerate}

\section{Batch Commit Tag Resolution}

The batch tag endpoint is one of the most impactful optimizations in the project.
The algorithm:

\begin{enumerate}
  \item rejects empty requests and requests larger than 100 commit IDs;
  \item resolves the repository ID and all branch and tag references;
  \item resolves tag references to their actual commit targets;
  \item builds a commit-to-tag map;
  \item creates a branch scan order with \code{main} first;
  \item scans branch histories recursively until all unresolved commits are found or all branches are exhausted;
  \item builds a result for each commit with \code{exactTags}, \code{tagAfter}, \code{tagBefore}, and \code{foundInBranch}.
\end{enumerate}

\screenshot[Batch tag enrichment result - commits grouped by repository with exact tags, before/after tags, and branch context]{Batch tag resolution results in the trace UI}

\section{Current UI Usage}

The batch-related frontend page groups discovered commits by repository and calls
\code{/commits/tags/batch} once per repository - not once per commit - giving
significantly better performance and a cleaner user experience.

\section{Conclusion}

The Commit module connects work-item and pull-request relationships back to
concrete code changes and repository references. It contains the most advanced
reference-resolution logic in the project.

%% ════════════════════════════════════════════════════════════════════════════
\chapter{Decision Module}

\section{Introduction}

The Decision module is a post-processing layer that receives already-collected
work-item and commit-tag data, then computes a decision-oriented aggregation
payload. It does not query Azure DevOps directly; it classifies input data into
stable result groups that can be consumed by upper layers.

\section{Implemented Endpoint}

\begin{center}
\code{POST /decision/aggregate}
\end{center}

\section{Input Contract}

The endpoint consumes \code{DecisionInputSetDto} with two arrays:

\begin{itemize}
  \item \code{workItems}: objects with \code{id} and optional \code{state};
  \item \code{repositories}: objects with \code{repositoryName} and commit entries containing \code{commitId}, \code{tagBefore}, \code{tagAfter}, and \code{exactTags}.
\end{itemize}

\noindent Validation is enforced through class-validator decorators
(\code{IsDefined}, \code{IsArray}, \code{ValidateNested}, \code{IsString}, and \code{IsNumber}).

\section{Aggregation Rules}

\code{DecisionAggregateService.aggregate(...)} applies deterministic rules:

\begin{enumerate}
  \item split work items into \code{doneWorkItems} and \code{otherWorkItems} using case-insensitive \code{state == "done"};
  \item for each repository, inspect each commit's \code{tagBefore}, \code{tagAfter}, and \code{exactTags};
  \item classify commits into three groups: before-and-after, only-before, and only-after;
  \item set \code{considerForPR=true} for only-before commits when the same repository contains at least one commit with an after/exact tag;
  \item emit repository-level arrays for each classification bucket.
\end{enumerate}

\section{Output Contract}

The result \code{DecisionAggregateResultDto} contains five top-level arrays:

\begin{itemize}
  \item \code{doneWorkItems}
  \item \code{otherWorkItems}
  \item \code{reposWithBeforeAndAfter}
  \item \code{reposWithOnlyBefore}
  \item \code{reposWithOnlyAfter}
\end{itemize}

\noindent This shape gives consumers a direct decision-ready view without
recomputing tag-position logic on the client side.

\section{Controller and Error Handling}

\code{DecisionController} is exposed under \code{/decision}, tagged in Swagger,
and protected by \code{MesGenericErrorFilter}. The controller logs input size and
result counts. Service failures are logged and converted to
\code{InternalServerErrorException} with a stable error message.

\screenshot[Decision aggregation flow - input work items and repository commit-tag references are classified into done/other work-item groups and before/after repository buckets]{Decision module aggregation flow}

\section{Conclusion}

The Decision module complements Work Item, Pull Request, and Commit modules by
transforming technical trace data into a structured aggregation contract tailored
for decision support.

%% ════════════════════════════════════════════════════════════════════════════
\chapter{Azure DevOps Integration and Configuration}

\section{Introduction}

This chapter describes how the backend connects to Azure DevOps, how secrets are
loaded through Dapr, and how the application is configured at runtime.

\section{AzureDevopsService}

The central integration point is \code{AzureDevopsService}. It implements
\code{OnModuleInit} so the Azure DevOps connection is established as soon as the
Nest module graph is ready.

\medskip
Its initialization flow:

\begin{enumerate}
  \item read the organization name, base URL, secret store name, and PAT secret name from configuration;
  \item create a \code{DaprClient};
  \item call \code{secret.get(...)} on the Dapr secret store;
  \item extract the PAT from the returned secret map;
  \item build the organization URL;
  \item create the \code{WebApi} connection using the PAT handler.
\end{enumerate}

\noindent Once connected, the service exposes typed accessors:
\code{getWorkItemTrackingApi()}, \code{getGitApi()}, \code{getCoreApi()}, and
\code{getBuildApi()}.

\screenshot[Dapr components folder structure showing the secret-store configuration used for PAT retrieval]{Dapr secret-store configuration}

\section{Configuration Modules}

Configuration is isolated through dedicated modules. \code{ConfigModule.forRoot(...)}
loads four separate configuration namespaces: application, URL, Azure DevOps, and
logger.

\section{Bootstrap Configuration}

The application bootstrap in \code{main.ts} configures:

\begin{itemize}
  \item \code{rawBody: true} for CloudEvents-compatible JSON parsing;
  \item a global validation pipe with \code{whitelist}, \code{forbidNonWhitelisted}, and \code{transform};
  \item explicit CORS for \code{http://localhost:5173} and \code{https://mesx-dev.app.corp};
  \item Swagger documentation at \code{/api};
  \item the service listening on port \texttt{3000}.
\end{itemize}

\screenshot[Swagger UI at \texttt{/api} - full endpoint list for the MES-X Dependency Management backend]{Swagger API documentation}

\section{Dapr Usage}

Dapr is used exclusively for secret retrieval in the current implementation.
The backend is future-ready for CloudEvents pub/sub, but no pub/sub routes are
active yet. The local development script \code{npm run start:dapr} starts the
service with the resources defined in \code{daprComponents/}.

\section{Frontend Configuration Boundary}

The frontend trace layer reads its backend base URL from
\code{URL\_DEPENDENCY\_PROTO\_BACKEND}. The value is normalized before use, and
the Axios client is configured with \code{withCredentials: true}. When the
backend returns \code{401} with a \code{loginUrl}, the UI redirects automatically.

\section{Conclusion}

The configuration and integration layers are critical enablers of the domain
modules. They keep sensitive values out of business services while maintaining a
clean, testable architecture.

%% ════════════════════════════════════════════════════════════════════════════
\chapter{Performance, Caching, and Reliability}

\section{Introduction}

This chapter explains the design choices that keep the service responsive and
predictable under realistic trace workloads.

\section{Batching Strategy}

Several operations are deliberately batched to reduce the number of API calls:

\begin{itemize}
  \item work-item fetching defaults to batches of 200 IDs;
  \item commit tag batch resolution rejects requests above 100 commit IDs;
  \item the frontend batch-related screen groups commit IDs by repository before calling the batch tag endpoint.
\end{itemize}

\section{Bounded Concurrency}

The project uses \code{p-limit} wherever unbounded parallelism could overload
Azure DevOps rate limits. Two clear examples:

\begin{itemize}
  \item PR expansion in the Work Item module with \code{p-limit(5)};
  \item PR commit expansion in the Commit module with \code{p-limit(5)}.
\end{itemize}

\section{Caching Strategy}

The backend contains two in-memory caches managed through \code{CacheManagerService}:

\begin{itemize}
  \item a work-item cache;
  \item a PR-to-work-item cache.
\end{itemize}

Both caches use FIFO-style eviction through \code{evictIfNeeded(...)} and are
capped at 1\,000 entries to prevent unbounded memory growth.

\section{Deduplication Strategy}

Deduplication is central to the correctness of every traversal. Maps and sets
are used at several levels:

\begin{itemize}
  \item root work-item IDs are normalized and deduplicated before traversal begins;
  \item discovered work items are tracked in a \code{Set<number>};
  \item visited PRs are tracked in a \code{Set<number>};
  \item fetched work items are stored in a map keyed by work-item ID;
  \item commits are deduplicated in a map keyed by commit SHA.
\end{itemize}

\section{Retry and Failure Handling}

The system includes several reliability measures:

\begin{itemize}
  \item each failed work-item batch is retried once before the request fails;
  \item PR expansion failures are logged and converted to empty sub-results when appropriate;
  \item missing work items produce explicit \code{NotFoundException} responses;
  \item invalid numeric identifiers are rejected early through validation helpers;
  \item the SSE stream broadcasts an error payload before closing clients when a failure occurs.
\end{itemize}

\section{SSE Stream Stability}

The streaming layer includes additional stability controls:

\begin{itemize}
  \item only one active stream is maintained at a time;
  \item new clients attach to the current stream if one already exists;
  \item all connected clients receive the same batch and completion payloads;
  \item the stream resets cleanly in both success and error cases.
\end{itemize}

\begin{infobox}[Current Architectural Gap]
The backend already exposes an aggregated batch related-items endpoint, but the
current UI batch results page still fans out single related-item requests per
selected root. This does not break functionality, but there is clear optimization
headroom for a future alignment pass.
\end{infobox}

\section{Conclusion}

Batching, bounded concurrency, cache management, deduplication, and defensive
error handling are the key reasons why the service remains usable even when the
trace graph becomes large and complex.

%% ════════════════════════════════════════════════════════════════════════════
\chapter{Testing, Tooling, and Delivery}

\section{Introduction}

This chapter summarizes the validation work, project scripts, and
delivery-related tooling surrounding both repositories.

\section{Backend Testing}

The backend uses Jest. The implemented tests cover the most critical parts of
the work-item and commit logic:

\begin{itemize}
  \item \code{work-item.controller.spec.ts}
  \item \code{work-item.service.spec.ts}
  \item \code{commit-history.service.spec.ts}
  \item \code{tag-selection.utils.spec.ts}
\end{itemize}

\noindent These tests validate success paths, error paths, deduplication behavior,
and the main-first tag resolution strategy.

\screenshot[Jest test output - all tests passing for work-item and commit history services]{Jest test results for the backend}

\section{Tooling and Scripts}

Table~\ref{tab:tooling-scripts} summarizes the most important scripts available
in both repositories.

\begin{table}[H]
\centering
\caption{Main tooling scripts}
\label{tab:tooling-scripts}
\rowcolors{2}{lightgray}{white}
\begin{tabularx}{\textwidth}{|p{2.5cm}|p{4.2cm}|X|}
\hline
\rowcolor{primary}
  \textcolor{white}{\textbf{Repo}} &
  \textcolor{white}{\textbf{Script}} &
  \textcolor{white}{\textbf{Purpose}} \\
\hline
Backend & \code{npm run build}       & Compile the NestJS service \\
Backend & \code{npm run lint}        & Lint TypeScript sources \\
Backend & \code{npm test}            & Run unit tests \\
Backend & \code{npm run test:e2e}    & Run end-to-end tests \\
Backend & \code{npm run start:dapr}  & Run backend with Dapr resources \\
Backend & \code{npm run typeorm:*}   & Generate, run, show, or revert migrations \\
UI      & \code{npm run dev}         & Start the Vite development server \\
UI      & \code{npm run build}       & Build the Vue trace UI \\
UI      & \code{npm run lint}        & Lint Vue and TypeScript sources \\
UI      & \code{npm run lintSass}    & Lint CSS, SCSS, and Vue style blocks \\
\hline
\end{tabularx}
\end{table}

\section{Frontend Integration Details}

The frontend already includes:

\begin{itemize}
  \item a router with the three trace routes;
  \item a Pinia store named \code{traceStore};
  \item a shared Axios API layer;
  \item composables such as \code{useBatchRelatedResults} and \code{useBatchStream};
  \item localized UI copy through \code{vue-i18n};
  \item Quasar-based components for tables, selection, status, and result cards.
\end{itemize}

\medskip\noindent The current batch user journey unfolds in three steps:

\begin{enumerate}
  \item start the SSE task stream and let work items load progressively;
  \item select the work items of interest from the loaded rows;
  \item navigate to the batch related-results page, where commits are grouped by
        repository and enriched with tag context through the batch tag endpoint.
\end{enumerate}

\screenshotwide[Full batch workflow - step 1: SSE task loading; step 2: work item selection; step 3: related results with tag enrichment]{Complete batch trace workflow (steps 1-3)}

\section{Delivery Assets}

Both repositories include delivery-oriented assets:

\begin{itemize}
  \item a \code{Dockerfile} for containerized deployment;
  \item \code{azure-pipelines.yml} for CI/CD automation;
  \item environment-variable support for runtime configuration;
  \item documented API contracts through Swagger and project documentation.
\end{itemize}

\section{Validation Reality}

The backend currently has the stronger automated validation surface. The frontend
has build and lint tooling in place, but does not yet include dedicated
trace-specific test coverage equivalent to the backend's Jest suite.

\section{Conclusion}

The project is not only implemented but also surrounded by practical tooling for
building, linting, testing, running locally, and preparing for deployment.

%% ════════════════════════════════════════════════════════════════════════════
\chapter{Final Result and Conclusion}

\section{Final Delivered Result}

By the end of the internship, the MES-X Dependency Management solution
comprises:

\begin{itemize}
  \item a NestJS backend exposing stable, read-only endpoints for work items,
        pull requests, commits, branches, and tags;
  \item a Vue\,3 trace UI with single-item trace, batch task streaming, and
        batch related-results visualization;
  \item shared contracts that include pull-request repository summaries through
        \code{visitedPullRequests};
  \item technical documentation aligned with the real implementation behavior.
\end{itemize}

\screenshotwide[Final application - all three trace screens: single item, batch stream, and batch results]{Overview of the delivered MES-X Dependency Management trace UI}

\section{What I Learned}

This internship strengthened several practical skills:

\begin{itemize}
  \item backend design with NestJS and TypeScript;
  \item DTO design and response normalization;
  \item integrating with Azure DevOps through specialized APIs;
  \item secure secret retrieval with Dapr;
  \item graph traversal, deduplication, and cache-aware recursion;
  \item batch-oriented UI integration with Vue, Pinia, and Axios;
  \item maintaining alignment between implementation, UI behavior, and documentation.
\end{itemize}

\section{Current Limitations and Future Improvements}

The current implementation is complete and usable. The following improvements
are planned for future iterations:

\begin{itemize}
  \item align the batch related-results UI with the existing backend batch endpoint,
        eliminating the per-root fan-out and improving overall performance;
  \item extend frontend automated test coverage around the three trace flows;
  \item continue refining API documentation and operational examples for new consumers.
\end{itemize}

\section{Conclusion}

This internship project was valuable because it addressed a real traceability
problem with concrete engineering work. It delivered structured backend modules,
a usable trace UI, and a set of algorithms that connect work-tracking artifacts
back to pull requests, commits, branches, and tags.

\medskip
Beyond the technical implementation, the project reinforced an important lesson:
keeping code, API contracts, UI behavior, and documentation in continuous
alignment is as important as the implementation itself. That discipline is what
makes a delivered solution maintainable by others long after the original
developer moves on.

%% ─── BIBLIOGRAPHY ────────────────────────────────────────────────────────────
\newpage
\addcontentsline{toc}{chapter}{Bibliography}
\renewcommand{\bibname}{Bibliography}
\begin{thebibliography}{99}

\bibitem{nestjs}
NestJS Team, \textit{NestJS Documentation}, 2024.
Available at: \url{https://docs.nestjs.com}

\bibitem{azdevops}
Microsoft, \textit{Azure DevOps REST API Reference}, 2024.
Available at: \url{https://learn.microsoft.com/en-us/rest/api/azure/devops}

\bibitem{dapr}
Dapr Authors, \textit{Dapr Documentation}, 2024.
Available at: \url{https://docs.dapr.io}

\bibitem{vue3}
Vue.js Team, \textit{Vue 3 Documentation}, 2024.
Available at: \url{https://vuejs.org/guide/introduction.html}

\bibitem{pinia}
Pinia Team, \textit{Pinia - The Vue Store that You Will Enjoy Using}, 2024.
Available at: \url{https://pinia.vuejs.org}

\end{thebibliography}

\end{document}